% from ega4-2-fr.tex
\oldpage[IV-2]{21}
morphisme ouvert. Mais cela résulte du critère (1.10.4) pour les morphismes
localement de présentation finie, les conditions \emph{b)}, \emph{b$'$)}, et
\emph{c)} de (1.10.4) n'étant autre que les conditions (i), (ii), (iii) de
(2.3.4).

\begin{corollary}[2.4.7]
\label{IV.2.4.7-fr}
Pour tout préschéma $Y$, le morphisme structural $\mathbf V_Y^n\to Y$ (où
\[
\mathbf V_Y^n
=Y\otimes_{\mathbf Z}\operatorname{Spec}
  \bigl(\mathbf Z[T_1,\ldots,T_n]\bigr),
\]
aussi noté $Y[T_1,\ldots,T_n]$) est universellement ouvert.
\end{corollary}

En effet, pour $Y=\operatorname{Spec}(A)$,
$Y[T_1,\ldots,T_n]=\operatorname{Spec}(A[T_1,\ldots,T_n])$, et
$A[T_1,\ldots,T_n]$ est une $A$-algèbre libre de présentation finie.

\begin{remark}[2.4.8]
\label{IV.2.4.8-fr}
\begin{enumerate}
\item[(i)] On notera qu'un morphisme $f:X\to Y$ fidèlement plat et
quasi-compact \emph{n'est pas nécessairement ouvert}, même lorsque $X$ et $Y$
sont noethériens. Prenons par exemple pour $Y$ une courbe algébrique réduite
irréductible (II, 7.4) de point générique $y$, et soit $X$ le préschéma somme
$Y\amalg\operatorname{Spec}(k(y))$, $f$ étant le morphisme canonique~; il est
clair que $f$ est plat et surjectif, donc fidèlement plat, et quasi-compact,
mais l'image par $f$ de la partie ouverte $\operatorname{Spec}(k(y))$ de $X$
est l'ensemble $\{y\}$, qui n'est pas ouvert dans $Y$ (II, 7.4.3).
\item[(ii)] Pour tout préschéma $X$, le morphisme canonique
$f:X_{\mathrm{red}}\to X$ est une immersion fermée et un
\emph{homéomorphisme universel} (2.4.4, (vi))~; mais lorsque $X$ est
localement noethérien, $f$ n'est plat que si $X$ est réduit, donc
$f=1_X$ (2.2.17).
\end{enumerate}
\end{remark}

\begin{proposition}[2.4.9]
\label{IV.2.4.9-fr}
Soit $Y$ un préschéma discret. Alors tout morphisme $f:X\to Y$ est
universellement ouvert.
\end{proposition}

La question étant locale sur $Y$ (2.4.3, (vii)), on peut se borner au cas où
l'espace sous-jacent $Y$ est réduit à un point~; remplaçant $f$ par
$f_{\mathrm{red}}$ (2.4.3, (vi)), on peut en outre supposer que $Y$ est le
spectre d'un \emph{corps} $k$~; d'autre part, pour tout changement de base
$Y'\to Y$, les ouverts de $X'=X\times_Y Y'$ images réciproques des ouverts
affines de $X$ recouvrent $X'$, donc on peut supposer
$X=\operatorname{Spec}(B)$ affine, $B$ étant une $k$-algèbre. Il s'agit donc
de prouver que pour toute $k$-algèbre $A'$, si l'on pose
$B'=B\otimes_k A'$, l'image par
$f':\operatorname{Spec}(B')\to\operatorname{Spec}(A')$ de tout ouvert de la
forme $U=D(t)$ ($t\in B'$) est ouverte dans $\operatorname{Spec}(A')$. Or,
$B$ est limite inductive de la famille filtrante croissante $(B_\alpha)$ de
ses sous-$k$-algèbres \emph{de type fini}, donc (le foncteur
$\varinjlim$ commutant au produit tensoriel), $B'$ est limite inductive des
$A'$-algèbres $B_\alpha\otimes_k A'=B'_\alpha$~; il existe $\alpha$ tel que
$t$ soit image dans $B'$ d'un élément $t_\alpha\in B'_\alpha$, donc $D(t)$
est l'image réciproque par le morphisme canonique
$u_\alpha:\operatorname{Spec}(B')\to\operatorname{Spec}(B'_\alpha)$ de
l'ouvert $U_\alpha=D(t_\alpha)$ (I, 1.2.2.2). Mais comme $k$ est un corps,
$B_\alpha$ est une $k$-algèbre de présentation finie et un $k$-module plat,
donc $B'_\alpha$ est une $A'$-algèbre de présentation finie et un $A'$-module
plat~; on conclut donc de (2.4.6) que l'image de $U_\alpha$ par
$f'_\alpha:\operatorname{Spec}(B'_\alpha)\to\operatorname{Spec}(A')$ est
\emph{ouverte} dans $\operatorname{Spec}(A')$. Tout revient par suite à voir
que $f'(U)=f'_\alpha(U_\alpha)$~; il est clair que
$f'(U)\subset f'_\alpha(U_\alpha)$ et il reste donc à voir que pour tout point
$y'\in f'_\alpha(U_\alpha)$, l'intersection
$V=U\cap f'^{-1}(y')$ n'est pas vide. Or on a
$V=u_\alpha^{-1}(V_\alpha)$, où
$V_\alpha=U_\alpha\cap f_\alpha'^{-1}(y')$~; autrement dit, $V_\alpha$ est
un ouvert non vide (par définition de $y'$) du préschéma
\[
\operatorname{Spec}(B'_\alpha\otimes_{A'}k(y'))
=\operatorname{Spec}(B_\alpha\otimes_k k(y'))
\]
et $V$ est son image réciproque par le morphisme
\[
v:\operatorname{Spec}(B\otimes_k k(y'))
\longrightarrow\operatorname{Spec}(B_\alpha\otimes_k k(y')).
\]
Comme $B_\alpha$ est une sous-algèbre de $B$ et que $k$ est un corps,
l'homomorphisme
$B_\alpha\otimes_k k(y')\to B\otimes_k k(y')$
